% Generated from validated Article Draft Result. Do not hand-edit; revise the structured draft instead.
\section{言語の外へ――Image GPTと離散visual token}
\label{pkg:media-autoregressive}
\noindent\textbf{2020年のgenerative mediaはdiffusionだけではない。Image GPTはpixel sequenceをautoregressiveにmodelingする方向を示し、Taming Transformersは高解像度画像をそのまま長大なsequenceとして扱う代わりに、CNNで学習した離散的なvisual vocabularyをTransformerで構成する設計を提示した。画像生成では、tokenとは何か、どの空間でsequence modelingするか自体が研究対象になった。}\autocite{src-03b656c5d291,src-d46ff01d5c8d}

Image GPTは、GPT系のautoregressive modelingを画像へ持ち込み、pixelをsequenceとして扱う研究方向を示した。languageとimageのdata structureは異なるが、前のtokenから次を予測するという一般的なsequence modelingの考え方をvisual domainへ適用した点に意味がある。2020年には、Transformer的な生成原理がtext以外でもどこまで成立するかが明示的な研究テーマになっていた。\autocite{src-d46ff01d5c8d}


Taming Transformersは、高解像度画像をpixel-levelの長いsequenceとして直接modelingすると計算量が大きくなる問題に対し、CNNを使ってcontext-richな離散visual vocabularyを学習し、そのcompositionをTransformerでmodelingする構成を示した。ここでは『画像を何にtokenizeするか』がarchitectureの中心になる。後年のlatent / discrete representation系へつながる設計空間を、2020年末の一次資料として確認できる。\autocite{src-03b656c5d291}


この系譜はDDPMなどのdiffusionと混同しない方がよい。Image GPTやTaming Transformersでは、離散的またはsequence-likeなrepresentationを順次modelingすることが中心であるのに対し、diffusion系ではnoiseを加えるforward processとdenoising / score estimationを通じてsampleを構成する。2020年のgenerative mediaは、一つの勝者へ収束する前に複数の生成原理が並立していた。\autocite{src-03b656c5d291,src-d46ff01d5c8d}


Image GPTからTaming Transformersまでを見ると、languageで強かったautoregressive sequence modelingをそのまま他modalitiesへ延長するだけでなく、画像側のrepresentationを変換してTransformerが扱いやすいsymbol列へ落とす方向も現れた。これは『生成AI=language model』という後年の単純化では捉えにくい2020年の分岐であり、翌年以降のmultimodal / generative-media研究を理解する前提になる。\autocite{src-03b656c5d291,src-d46ff01d5c8d}


\begin{claimboundary}
Image GPTとTaming Transformersの生成品質やbenchmark評価はauthor / project側の報告である。本稿はautoregressive image modelingと離散visual representationという設計上の分岐を一次資料から確認しているが、画像品質やstate-of-the-art主張を独立再現したものではない。\autocite{src-03b656c5d291,src-d46ff01d5c8d}
\end{claimboundary}
